\section{Proof Sketch}
\label{sec:proof-sketch}

The scalar identity $\sigma(r)-\sigma(-r)=r$ makes the antisymmetric part of
each realized depth-two network exactly linear.  Summing over the prescribed
latter-half iterates gives
$A_\omega(x)=\langle v_\omega,x\rangle$; see
Proposition~\ref{prop:step-001-aggregate}.

Oddness and the fixed tie rule imply a pointwise pair-error comparison.  If
$G_\omega$ correctly labels both $x$ and $-x$, one of the two correctness
conditions is strict even when the other score is zero, so
$G_\omega(x)-G_\omega(-x)$ has the strict sign of $h(x)$.  Integrating this
comparison gives Proposition~\ref{prop:step-002-risk-transfer},
\[
\mathcal L_{\mathcal D,h}(A_\omega)
\leq2\mathcal L_{\mathcal D^{\mathrm{sym}},h}(G_\omega).
\]
The universal source premise is then invoked on the legal distribution
$\mathcal D^{\mathrm{sym}}$.  A finite-range expectation argument extracts,
for every original pair $(\mathcal D,h)$, a deterministic $v$ satisfying
\[
\mathcal L_{\mathcal D,h}(x\mapsto\langle v,x\rangle)
\leq2\varepsilon.
\]

For a fixed target, $Q_h$ contains exactly the representative labeled
$-s_0$ from each antipodal pair.  Exact homogeneous representation is
equivalent to feasibility of the strict system
$h(q)\langle w,q\rangle>0$ on $Q_h$.  If this system were infeasible, the
closest point of the convex hull of the signed vectors $h(q)q$ would be zero.
Affine-dependence pruning reduces the resulting positive convex certificate
to at most $n+1$ representatives.  On the uniform distribution over those
representatives, every homogeneous score makes at least one tie-resolved
error, so its risk is at least $1/(n+1)$.

Applying the deterministic $2\varepsilon$ upper bound to that same witness
distribution contradicts Assumption~\ref{assump:high-accuracy}.  Thus the
strict system is feasible for every target and the identity map represents
the class exactly.  The Dirac law on the identity map then succeeds with
probability one.  Finally, every deterministic representation induces a
confident one by taking a point mass, and
$n\leq m(n+1)=S\leq TS$, which yields the asserted chain.
